\documentclass[UTF8,a4paper,12pt]{ctexrep}
\usepackage{amsmath}
\usepackage{array}
\usepackage{geometry}
\geometry{left=2.8cm,right=2.8cm,top=2.6cm,bottom=2.6cm}
\title{爬壁机器人船体表面焊缝跟踪控制与运动优化研究}
\author{}
\date{}
\begin{document}
\maketitle
\tableofcontents

\chapter*{爬壁机器人船体表面焊缝跟踪控制与运动优化研究}

\section{目录}

第一章 绪论 1.1 研究背景与意义 1.2 国内外研究现状 1.3 本文研究目标与主要内容 1.4 本文技术路线 1.5 本文组织结构 1.6 本章小结\\

第二章 视觉感知与目标位置表征方法 2.1 焊缝跟踪任务中的视觉感知需求 2.2 图像输入与视觉检测流程 2.3 基于目标检测的焊缝区域识别方法 2.4 焊缝目标中心位置表征方法 2.5 图像偏差构建与控制输入转换 2.6 本章小结\\

第三章 运动控制与焊缝跟踪方法 3.1 焊缝跟踪控制任务分析 3.2 麦克纳姆轮全向底盘速度接口与运动学模型 3.2.1 底盘执行层速度分配与电机闭环实现 3.3 基于视觉偏差的跟踪误差构建 3.4 面向视觉闭环的运动优化控制策略 3.5 感知—控制闭环机制 3.6 异常处理与安全停止逻辑 3.7 本章小结\\

第四章 实机验证与仿真分析 4.1 系统实现环境与总体流程 4.2 视觉识别功能验证 4.3 焊缝跟踪闭环运行验证 4.4 Gazebo 仿真场景与模型支撑分析 4.5 不同焊缝场景下的运行现象分析 4.6 实验结果与不足 4.7 本章小结\\

第五章 总结与展望 5.1 全文工作总结 5.2 方法与系统实现总结 5.3 当前工作的不足 5.4 后续研究展望\\

\section{论文题目}

《爬壁机器人船体表面焊缝跟踪控制与运动优化研究》\\

\section{摘要}

焊缝跟踪是爬壁机器人开展船体表面焊缝检测、辅助作业和路径跟随任务中的重要基础环节。为了使机器人能够根据视觉信息获得焊缝目标位置并生成运动控制量，本文围绕“视觉感知、目标位置表征、偏差构建、运动控制和闭环验证”这一技术路线，设计并实现了一种基于 YOLOv5 目标检测的焊缝跟踪方法。系统以相机图像或测试视频帧作为输入，通过目标检测模型获得焊缝区域检测框，并进一步提取检测框横向中心作为焊缝目标位置表征。在此基础上，以图像中心作为期望参考位置，建立目标中心与图像中心之间的横向偏差关系，将视觉识别结果转换为运动控制输入。\\

在运动控制方面，本文面向麦克纳姆轮全向移动底盘的速度接口，采用基于视觉偏差的速度调节方法。系统根据焊缝目标相对图像中心的偏移方向和偏移程度生成角速度控制量，并围绕视觉闭环跟踪需求进行运动优化设计：首先通过归一化误差降低图像尺度变化对控制量的影响，然后在比例控制基础上加入小增益积分项，并结合积分分离、积分限幅和线速度-角速度耦合约束，使机器人在视觉目标引导下完成对焊缝目标的局部跟踪运动。当视觉检测无效或中心点输入超时时，系统输出零速度指令，以降低目标丢失情况下继续运动的风险。系统实现基于 ROS1/catkin 组织，包含图像输入、视觉检测、位置表征、速度控制、执行接口和 Gazebo 仿真支撑等环节，形成了从图像感知到运动输出的闭环处理流程。\\

系统实现结果表明，所构建的焊缝跟踪系统已具备图像输入、目标检测、结果可视化、中心位置输出、速度指令生成、底盘接口和仿真场景支撑等功能基础。由于目前尚未形成完整的定量实验结果、模型训练日志、检测精度指标和实机跟踪误差数据，本文实验部分以可行性验证流程和待补充实验材料说明为主，不对识别率、平均误差、运行速度等指标作确定性结论。研究结果表明，基于检测框中心的视觉位置表征能够为焊缝跟踪控制提供简洁输入，基于图像中心偏差的速度调节方法能够支撑视觉感知与机器人运动控制之间的连续连接，为后续实机实验、参数优化和鲁棒性提升提供了基础。\\

\section{关键词}

爬壁机器人；船体表面；焊缝跟踪；目标检测；YOLOv5；视觉偏差控制；运动优化；ROS1\\

\section{摘要写作边界说明}

当前摘要没有写入具体检测准确率、mAP、FPS、平均误差、工业现场测试和完整实机部署结论。若后续用户补充可靠实验数据，可在摘要最后增加定量结果描述。\\

\chapter*{第一章 绪论}

\section{1.1 研究背景与意义}

船体结构由大量钢板、型材和舱段焊接形成，焊缝质量直接关系到船体结构强度、密封性能和服役安全。船舶建造、维护和检测过程中，船体表面作业区域通常具有面积大、位置分散、曲面变化明显和作业环境复杂等特点，人工检测和人工跟踪不仅效率较低，而且容易受到高处作业、视野受限、操作疲劳和环境干扰的影响。张永杰等在船用机器人爬壁平台综述中指出，船用爬壁机器人区别于普通地面机器人之处在于其同时具备吸附和移动能力，并可搭载清洗、焊接、图像识别等装置完成船体表面作业[1]。刘冰杰围绕船体除锈维护任务设计了磁吸附爬壁机器人样机，并从吸附稳定性、负载能力和越障能力等方面验证船体维护场景对平台能力的要求[2]。因此，利用爬壁机器人代替或辅助人工完成船体表面焊缝相关作业，是船舶制造与维护装备智能化的重要发展方向。\\

焊缝跟踪并不只是机器人移动平台问题，还包含视觉感知、目标位置表征和运动控制之间的连续转换。对于船体表面焊缝跟踪任务而言，机器人首先需要在复杂表面上获得稳定的焊缝视觉信息，然后需要把图像中的焊缝区域转化为可用于控制的位置信息，最后根据目标位置偏差生成速度控制量。Zhang 等面向船舶焊缝磁粉检测任务设计了轮式爬壁机器人，通过永久磁吸附、柔顺检测机构和自动检测流程提高船舶焊缝检测效率，说明焊缝作业机器人需要同时考虑平台运动、检测执行和流程控制之间的配合[3]。这类研究表明，面向船体表面的焊缝跟踪控制方法具有明确的工程需求和研究价值。\\

从技术链路看，焊缝跟踪不是单一的图像识别问题，而是视觉感知与运动控制之间的耦合问题。视觉系统获得的是二维图像，图像中的焊缝目标可能表现为线状区域、纹理区域或具有一定宽度的目标区域；运动控制系统需要的则是线速度、角速度等可执行控制量。若系统只能在图像中识别出焊缝，而不能进一步得到焊缝相对于机器人运动方向的偏差信息，控制系统仍然无法判断机器人应保持前进、转向修正还是停止。因此，焊缝跟踪任务的关键在于建立从图像目标、位置表征、偏差构建到运动控制输出之间的转换关系。\\

近年来，基于机器视觉和深度学习的焊缝识别方法不断发展，传统阈值分割、边缘检测、霍夫变换、激光轮廓传感、目标检测、语义分割和多源视觉融合等方法均被用于焊缝识别或路径跟踪。Cui 等针对船体曲面和障碍跨越问题设计了轮腿复合磁吸附爬壁机器人，体现了船体表面运动平台对复杂曲面适应性的需求[4]。在焊缝视觉识别方面，Li 等将 Mask R-CNN 用于爬壁检测机器人的焊缝识别与跟踪，通过实例分割结果提取焊缝路径信息[12]；Wang 等采用 DeepLabv3+ 语义分割模型识别焊缝，并结合轻量化网络结构与路径拟合方法提高系统部署可行性[13]；Li 等将改进 YOLOv5 与 RealSense 深度相机结合，用于机器人焊接引导中的焊缝坡口定位[15]；Wang 等探索二维图像与三维点云融合的焊缝检测方法，以增强焊缝定位结果的可靠性和可解释性[17]。这些成果为本文从检测框目标位置出发构建焊缝跟踪控制方法提供了研究背景。\\

本文题目为“爬壁机器人船体表面焊缝跟踪控制与运动优化研究”。需要说明的是，本文的研究重点不是爬壁机器人机械结构设计、磁吸附机构设计或船体曲面力学分析，而是面向爬壁机器人船体表面焊缝跟踪任务，研究上层视觉感知、目标位置表征、视觉偏差构建、速度控制输出和安全停止逻辑。当前系统以相机图像或测试视频图像为输入，通过 YOLOv5 推理得到焊缝目标检测框，提取检测框横向中心和图像宽度，并以图像中心作为期望参考构建横向偏差；控制部分根据偏差生成角速度和线速度控制量，并在检测无效或输入超时时输出零速度。由此，系统形成了从视觉感知到运动控制输出的连续处理链路。\\

本文所使用的平台在代码材料中主要体现为麦克纳姆轮全向移动底盘及其速度接口，相关材料尚未提供真实船体爬壁吸附结构、船体表面实测环境和完整实机定量实验结果。因此，本文将“爬壁机器人”作为面向船体表面焊缝作业的应用对象和研究背景，将重点放在焊缝跟踪控制和运动优化方法上。参考文献中涉及的磁吸附结构、轮腿机构、Mask R-CNN、DeepLabv3+、三维视觉、Stanley 控制、模糊 PID 和多传感器融合等内容，只作为国内外研究现状与背景对比，不写成本文已经实现的内容。\\

\section{1.2 国内外研究现状}

\subsection{1.2.1 船体表面爬壁机器人研究现状}

船体表面爬壁机器人研究首先关注机器人如何在钢制壁面、曲面和焊缝突起附近保持可靠吸附和稳定运动。张永杰等按照吸附原理和运动特点，将船用爬壁平台归纳为磁吸附、负压吸附、仿生吸附等类型，并从机械结构、运动性能和适用范围方面比较了不同平台的特点；该综述还指出，船体表面常存在曲率变化、焊接突起、壁面变形和腐蚀等因素，单纯依赖刚性结构难以适应船体复杂作业环境[1]。这一研究为船体表面爬壁平台选型提供了系统参考，也说明焊缝跟踪任务需要考虑船体表面环境对机器人运动稳定性的影响。\\

面向具体船体维护作业，刘冰杰设计了携带高压射流除锈装置的船体爬壁机器人，围绕磁吸附模块、齿轮传动模块、驱动模块和除锈清洗模块进行结构设计，并通过力学分析、磁场仿真和样机实验验证直行、转弯、负载与越障性能[2]。该研究的重点在于船体表面维护机器人的结构设计和吸附能力验证，能够说明船体场景下机器人平台需要具备较强承载和越障能力；但其研究重心并不是焊缝视觉跟踪控制。\\

Zhang 等针对船舶焊缝磁粉检测任务，提出一种基于永久磁吸附的轮式爬壁机器人，并设计串并联柔顺检测机构，使磁轭和检测区域能够更好贴合船体曲面；该研究还建立统一力学模型分析不同姿态下的吸附与运动条件，并提出符合磁粉检测流程的自动检测控制流程[3]。该文献说明，在船舶焊缝检测场景中，平台运动稳定性、检测机构贴合性和作业流程控制需要共同考虑。与该类研究相比，本文不设计磁粉检测机构，而是把船体焊缝作业作为应用背景，重点研究视觉目标位置如何进入运动控制。\\

Cui 等提出用于船体外板的磁吸附爬壁机器人，采用仿生启发的轮腿复合移动结构，在保持轮式运动效率的同时改善越障和曲面适应能力；文中还对机器人运动学模型、路径跟踪方法和不同预瞄距离对路径跟踪的影响进行了分析[4]。该研究强调船体曲面和障碍对爬壁机器人运动控制的影响，为焊缝跟踪中的运动稳定性提供了参考。但本文当前系统并未实现轮腿机构或纯跟踪路径规划，只借鉴其“复杂船体表面需要稳定运动控制”的问题意识。\\

除焊缝检测外，船体表面损伤检测和焊缝打磨机器人也从不同角度推动了船体爬壁平台的发展。余加林面向船体表面损伤探测设计四轮磁吸附爬壁机器人，并融合曲率传感和漏磁检测用于船体损伤识别，体现了船体检测任务对多传感信息采集和稳定壁面运动的需求[5]。江来针对船体表面焊缝打磨任务，研究了轮式磁吸附爬壁机器人、深度相机点云识别、焊缝打磨路径规划和控制系统设计，说明焊缝相关作业不仅需要到达焊缝区域，还需要对焊缝几何位置进行识别和路径生成[6]。这些研究与本文的共同点在于均面向船体表面焊缝或损伤作业，不同之处在于本文不开展打磨路径规划和三维点云处理。\\

曹爽爽围绕磁吸附爬壁机器人开展结构设计、磁吸附模块仿真、履带机构设计和样机性能实验，分析了机器人在不同壁面姿态下防滑移、防倾覆所需的磁吸附力，并通过直行、转向、负载和越障实验验证样机性能[7]。该研究体现了爬壁机器人本体设计中常见的吸附安全和力学约束问题。本文在绪论中引用该类文献，是为了说明爬壁机器人平台研究已经形成较多机械结构和吸附分析成果；本文后续章节不把这些内容写成本系统已经完成的机械结构设计。\\

\subsection{1.2.2 焊缝识别与跟踪技术研究现状}

焊缝识别与跟踪技术早期多依赖传统图像处理、激光视觉或轮廓传感器。Zhang 等设计的焊缝检测机器人采用四轮底盘、柔性检测头和基于 YOLOv5 的焊缝目标识别算法，同时使用改进 TEB 算法进行检测路线规划；该研究更侧重自主到达检测位置和识别焊缝位置，文中也指出其系统主要定位焊缝并辅助检查人员观察，而没有进一步完成复杂缺陷分类[8]。这一研究说明，视觉识别可以作为焊缝检测机器人到达与定位的关键环节，但完整焊缝检测系统还需要路径规划、运动控制和人工或自动检测流程配合。\\

王季面向大型压力容器焊缝检测爬壁机器人，采用图像灰度化、高斯滤波、Canny 边缘检测、概率霍夫直线检测和最小二乘拟合等方法提取焊缝中心位置，并在此基础上结合积分分离式 PID 控制完成机器人焊缝自主跟踪[9]。该研究的特点是图像处理流程清晰、控制接口直接，适合焊缝对比度较明显的场景；其局限在于焊缝图像质量、光照变化和背景干扰会影响边缘与直线特征提取效果。\\

陈晨等针对大型承压设备顶部和侧壁焊缝的 PAUT、TOFD 检测需求，设计了基于激光轮廓传感器的焊缝智能跟踪爬壁检测机器人。该方法利用轮廓传感器获取焊缝余高轮廓，通过多尺度直线拟合法拾取轮廓数据突变点，再根据焊缝中心偏差调整爬行器左右驱动模块线速度，实现对直焊缝和弯曲焊缝的跟踪[10]。该研究说明，激光轮廓传感器能够利用焊缝几何突变获得较稳定的焊缝中心信息，但系统需要额外的轮廓传感硬件和相应标定，硬件复杂度高于单目图像目标检测方案。\\

宋奇恒面向船体合拢焊缝外观缺陷检测，提出涵盖图像采集、焊缝跟踪和缺陷检测的系统方案，并使用磁吸附爬壁机器人作为载体；其焊缝跟踪部分采用改进旋转框检测网络识别焊缝位姿，再根据检测输出设计机器人纠偏策略[11]。该研究对象与本文同属船体焊缝场景，说明船体焊缝在实际船坞环境中存在复杂光照、焊缝尺度大和微小缺陷难检测等问题。本文不采用 YOLOv8 旋转框检测，也不写入船坞现场定量指标，而是将其作为船体焊缝视觉跟踪研究的近年参考。\\

综合传统图像处理、激光轮廓传感和旋转框检测研究可以看出，焊缝识别与跟踪的核心并不只是“找到焊缝”，还包括如何把焊缝的图像或轮廓结果转化为机器人控制可使用的偏差信息。传统方法计算量较低、控制接口直观，但通常依赖稳定光照、明显边缘或专用传感器；深度学习方法适应性更强，但需要训练数据、推理部署和输出解释机制支撑。本文选择检测框中心作为目标位置表征，正是为了在实现复杂度和控制可用性之间取得折中。\\

\subsection{1.2.3 基于深度学习的焊缝视觉检测研究现状}

随着深度学习的发展，焊缝视觉检测逐渐从人工特征提取转向学习型特征表达。Li 等将 Mask R-CNN 用于爬壁检测机器人焊缝识别与跟踪，机器人采用麦克纳姆轮和永久磁吸附结构，通过实例分割获得焊缝掩膜，再结合图像处理和霍夫变换提取焊缝路径，并根据漂移角和偏移距离调整机器人运动[12]。该研究展示了实例分割方法在焊缝路径提取中的能力，但 Mask R-CNN 需要像素级掩膜输出和较高计算资源。本文并未实现 Mask R-CNN 或焊缝中心线拟合，只将其作为深度学习焊缝识别研究背景。\\

Wang 等设计了集成焊缝跟踪与检测功能的壁面机器人，使用 DeepLabv3+ 语义分割模型识别焊缝，并通过 MobileNetv2、注意力机制和深度可分离空洞卷积对模型进行轻量化改进，随后基于分割结果采用最小二乘方法拟合焊缝路径[13]。该研究说明，语义分割能够获得比检测框更细致的焊缝区域信息，适合需要路径拟合或精细检测的任务；但这类方法需要更复杂的数据标注和模型部署。本文当前系统只使用目标检测框中心，不写成已经实现语义分割或路径拟合。\\

Du 等提出基于机器视觉的自动焊缝检测爬壁机器人系统，将爬壁平台与视觉检测结合用于焊缝自动检测，体现了机器视觉在爬壁机器人焊缝作业中的系统集成趋势[14]。与该类系统相比，本文的工作更聚焦于检测框中心、图像中心和速度控制之间的接口关系，不对焊缝检测精度、缺陷识别能力或工业现场运行效果作超出材料依据的结论。\\

Li 等针对机器人焊接引导任务，将改进 YOLOv5 与 RealSense 深度相机结合，在 YOLOv5 中引入坐标注意力模块以提高坡口检测能力，再利用深度信息计算焊缝坡口在机器人坐标系中的空间位置[15]。该研究说明，YOLOv5 检测框中心可以与深度相机信息结合，用于从二维图像位置扩展到三维引导位置。本文同样采用 YOLOv5 检测框作为视觉前端输出，但当前材料未体现 RealSense 深度相机、手眼标定或三维坐标转换，因此只研究二维图像中心偏差控制。\\

Zhu 等面向钢结构焊缝识别任务，在 YOLOv5 中引入坐标注意力机制并构建 C3CA 模块，以增强模型对焊缝空间位置的感知能力，在保证检测精度的同时降低部分参数量和计算量[16]。该研究体现了 YOLO 系列模型在焊缝目标定位中的轻量化改进趋势。本文没有改造 YOLOv5 网络结构，而是使用现有 YOLOv5 推理链路完成焊缝目标框输出，并进一步研究该输出如何服务于运动控制。\\

除二维目标检测和分割外，多源视觉融合也是焊缝检测的重要方向。Wang 等将二维图像处理与三维点云分析结合，采用改进 Faster R-CNN 和正交平面交线提取算法提升焊缝检测可靠性，并通过可解释特征图和透明候选框生成过程增强模型输出的可解释性[17]。Pham 等则利用激光线和固定视角 RGB 相机构建低成本焊缝跟踪模型，通过相机标定、颜色阈值分割和随机森林、梯度提升回归等方法实现图像坐标到实际坐标的转换[18]。这些研究说明，三维视觉和机器学习坐标映射可以提升焊缝定位的空间表达能力，但也需要标定、额外传感器或复杂数据处理。本文不涉及三维点云与世界坐标转换，而采用更简单的图像横向中心偏差作为控制输入。\\

\subsection{1.2.4 焊缝跟踪控制与运动优化研究现状}

焊缝跟踪控制的目标是将焊缝识别结果转化为机器人运动控制量，使机器人沿焊缝方向稳定运动。邢正等针对球罐爬壁机器人使用单激光传感器难以双向稳定跟踪球罐焊缝的问题，结合球面运动学特性进行误差投影与补偿，并在 Stanley 模型基础上引入模糊自适应控制动态调整速度增益系数[19]。该研究说明，对于曲面焊缝跟踪任务，单纯的平面误差控制可能不足，需要结合曲面运动学和速度参数自适应机制。本文不实现 Stanley 控制或球面误差补偿，而是借鉴其将传感器误差转换为运动控制量的思路。\\

在可为焊缝跟踪运动优化提供参考的移动机器人控制研究中，模糊控制和参数优化方法常用于改善运动稳定性。邓弘文等针对四轮差速机器人运动控制问题，将优化算法与模糊控制结合，用于改善机器人运动过程中的稳定性和控制性能[20]。张树晓针对光伏板清洁机器人防偏摆运动控制问题，采用 PSO 优化模糊 PID 参数，以提高机器人在作业过程中的抗偏摆能力和运行平稳性[21]。这些研究并非直接面向船体焊缝跟踪，但能够说明速度控制参数、误差反馈和抗扰优化对机器人作业稳定性具有参考价值。本文当前系统采用小增益积分修正、积分限幅和线速度-角速度耦合约束，属于更简化的运动优化策略。\\

邱东围绕室内寻迹机器人运动控制优化，结合 Frenet 坐标系描述路径跟踪误差，并从视觉传感和轨迹优化角度讨论寻迹控制方法[22]。该研究说明，在路径跟踪问题中，将视觉信息转换到统一的误差描述空间是控制设计的重要环节。本文不采用 Frenet 坐标系，也不进行全局轨迹优化，而是直接以图像中心作为期望参考、以检测框中心作为当前目标位置，构建横向视觉偏差。\\

总体来看，已有焊缝跟踪控制研究已经形成了从传统 PID、模糊控制、Stanley 模型到视觉伺服和路径规划等多种方法。复杂方法通常需要较完整的路径几何信息、车辆状态反馈、曲面模型或多传感器定位数据，适合高精度或复杂工业场景；而本文当前系统依据代码材料可确认的事实，主要具备 YOLOv5 检测框输出、图像中心偏差计算、速度指令生成和无检测停止能力。因此，本文将运动优化限定在视觉偏差闭环内部，重点讨论归一化偏差、小增益积分修正、线速度-角速度耦合约束和安全停止机制，而不把未实现的路径规划、三维定位或高级控制算法写成本系统成果。\\

综合以上研究现状可以看出，船体表面焊缝跟踪涉及爬壁平台、视觉感知、目标位置表征和运动控制等多个层面。已有文献在机械结构、吸附设计、深度学习分割、三维视觉和高级控制方面提供了丰富参考，但本科课题在时间、实验条件和代码基础方面存在现实边界。本文的切入点是在不展开爬壁机械本体设计的前提下，面向船体表面焊缝作业背景，构建一条简洁的视觉感知到运动控制链路：由 YOLOv5 得到焊缝检测框，提取检测框横向中心，计算其与图像中心的偏差，再生成速度控制量并在检测无效时停止。该路线既与焊缝跟踪主题相吻合，也符合当前系统实现和论文事实边界。\\

\section{1.3 本文研究目标与主要内容}

本文研究目标是面向爬壁机器人船体表面焊缝跟踪任务，构建一种由视觉目标检测结果驱动运动控制输出的方法体系。该体系应能够从输入图像中获得焊缝目标区域，提取目标中心位置，构建目标中心与图像中心之间的偏差，并根据偏差信息生成速度控制量。\\

围绕上述目标，本文主要研究内容包括以下四个方面。\\

第一，研究焊缝目标视觉检测流程。系统通过统一图像输入获取相机或测试视频图像，经过图像预处理后调用 YOLOv5 推理方法获得焊缝目标检测框。该部分重点关注检测框如何作为焊缝区域的位置表达，而不展开未确认的训练过程和模型精度。\\

第二，研究焊缝目标位置表征方法。系统从检测框左右边界计算目标横向中心，并结合图像宽度构造图像中心参考。通过目标中心和图像中心之间的关系，建立横向偏差，为运动控制提供输入。\\

第三，研究基于视觉偏差的运动控制方法。系统根据目标横向偏差生成角速度控制量，并根据偏差大小约束线速度。当目标接近图像中心时，机器人保持前进；当目标偏离图像中心时，机器人根据偏差方向转向；当目标检测无效或输入超时时，系统输出零速度。\\

第四，研究系统实现与验证流程。系统基于 ROS1/catkin 组织，包含图像输入、视觉检测、位置表征、控制输出、底盘接口和 Gazebo 仿真支撑。由于目前尚未形成船体表面爬壁实机实验和完整定量结果，本文验证部分以可行性验证、运行现象分析和待补充材料说明为主。\\

从工作边界看，本文使用现有移动底盘速度接口、ROS1 通信框架和 YOLOv5 推理工具作为基础支撑，但论文工作并不止于工具调用。本文完成的核心工作包括：将目标检测框输出转化为控制可用的位置表征，建立目标中心与图像中心之间的归一化偏差关系，构建检测有效性与控制安全状态之间的对应关系，并设计由角速度修正、线速度约束、积分限制和超时停止组成的运动优化控制策略。上述工作使机器人能够围绕焊缝目标形成专门的感知—控制闭环，而不是仅依赖人工遥控或平台原有运动能力。\\

\section{1.4 本文技术路线}

本文技术路线围绕视觉感知到运动控制的连续转换展开，可抽象为：\\

\[
I_t \rightarrow D_t \rightarrow p_t \rightarrow e_t \rightarrow u_t \rightarrow q_{t+1}
\]
其中，$I_t$ 表示 $t$ 时刻输入图像，$D_t$ 表示视觉检测结果，$p_t$ 表示目标位置表征，$e_t$ 表示偏差信息，$u_t$ 表示控制输出，$q_{t+1}$ 表示机器人下一时刻运动状态。该公式用于概括感知—控制闭环中的信息传递关系，是对本文系统流程的论文式抽象。\\

具体而言，系统首先获取相机或测试视频图像；然后通过目标检测方法得到焊缝目标检测框；接着提取检测框横向中心，并与图像中心比较得到横向偏差；随后根据偏差生成角速度和线速度；最后由底盘接口或仿真插件执行速度指令。若视觉检测无效，则系统进入零速度停止状态。该技术路线体现了本文的核心思想：不是简单连接若干工程模块，而是围绕焊缝目标位置如何驱动机器人运动形成完整处理链。\\

\section{1.5 本文组织结构}

本文共分为五章。\\

第一章为绪论，主要说明爬壁机器人船体表面焊缝跟踪任务的研究背景与意义，分析船体表面爬壁机器人、焊缝识别与跟踪、深度学习焊缝视觉检测以及焊缝跟踪控制与运动优化相关研究现状，并给出本文研究目标和技术路线。\\

第二章为视觉感知与目标位置表征方法，重点说明图像输入、YOLOv5 检测流程、检测框中心提取、图像中心参考和视觉偏差构建方法。\\

第三章为运动控制与焊缝跟踪方法，围绕视觉偏差如何转化为速度控制量展开，说明麦克纳姆轮全向底盘速度接口、车体速度层运动学关系、角速度生成、线速度约束、运动优化控制策略、闭环机制和安全停止逻辑。\\

第四章为实机验证与仿真分析，说明系统验证目的、视觉识别功能验证、闭环运行验证、Gazebo 仿真支撑和当前实验材料不足。\\

第五章为总结与展望，总结全文工作，分析当前系统在训练数据、定量评价、实机验证和控制优化方面的不足，并提出后续研究方向。\\

\section{1.6 本章小结}

本章从爬壁机器人船体表面焊缝跟踪任务需求出发，说明了视觉感知与运动控制之间建立连续关系的必要性。随后参考代表性文献，对船体表面爬壁机器人、焊缝识别与跟踪技术、基于深度学习的焊缝视觉检测方法以及焊缝跟踪控制与运动优化研究进行了综述。已有研究为本文提供了平台背景、视觉识别方法和控制优化思路，但本文不展开爬壁机械结构、磁吸附机构、三维视觉和高级路径跟踪算法，而是将系统定位为面向船体表面焊缝作业背景、基于 YOLOv5 目标检测和图像中心偏差控制的焊缝跟踪系统。后续章节将在明确事实边界的基础上，对视觉感知方法、运动控制方法和验证分析进行展开。\\

\section{本章依据说明}

本章依据 `current\_thesis\_code\_handoff/code\_based\_thesis\_handoff.md` 中“项目总体定位”“整体技术路线”“核心方法与原理”“系统实现逻辑”和“明确不能写的内容”。本章文献引用主要用于支撑船体表面爬壁机器人、焊缝检测、焊缝视觉识别和跟踪控制研究现状；相关文献中的机械结构设计、吸附设计、语义分割、三维视觉、Stanley 控制、模糊 PID 和定量实验结果均作为研究背景，不写成本文已经实现的内容。\\

\section{参考文献}

[1] 张永杰, 张诺, 杨振, 等. 船用机器人爬壁平台研究进展[J]. 应用力学学报, 2024, 41(1): 25-35. [2] 刘冰杰. 船体爬壁机器人设计与研究[D]. 燕山大学, 2024. [3] Zhang X, Zhang M, Jiao S, Sun L, Li M. Design and Optimization of the Wall Climbing Robot for Magnetic Particle Detection of Ship Welds[J]. Journal of Marine Science and Engineering, 2024, 12(4): 610. [4] Cui S, Pei X, Song H, Dai P. Design and Motion Analysis of a Magnetic Climbing Robot Applied to Ship Shell Plate[J]. Machines, 2022, 10(8): 632. [5] 余加林. 船体表面损伤探测爬壁机器人设计及研制[D]. 哈尔滨工业大学, 2025. [6] 江来. 船体表面焊缝打磨机器人设计及关键技术研究[D]. 江苏科技大学, 2024. [7] 曹爽爽. 一种磁吸附爬壁机器人的设计与分析[D]. 燕山大学, 2024. [8] Zhang P, Wang J, Zhang F, Xu P, Li L, Li B. Design and analysis of welding inspection robot[J]. Scientific Reports, 2022, 12: 22651. [9] 王季. 大型压力容器焊缝检测爬壁机器人设计与研究[D]. 安徽工业大学, 2022. [10] 陈晨, 江志铭, 马源旺, 等. 基于轮廓传感器的焊缝智能跟踪爬壁检测机器人[J]. 南昌航空大学学报: 自然科学版, 2025, 39(3): 92-100. [11] 宋奇恒. 船体合拢焊缝外观缺陷检测系统研究与验证[D]. 机械科学研究总院, 2025. [12] Li J, Li B, Dong L, Wang X, Tian M. Weld Seam Identification and Tracking of Inspection Robot Based on Deep Learning Network[J]. Drones, 2022, 6(8): 216. [13] Wang J, Huang L, Yao J, Liu M, Du Y, Zhao M, et al. Weld Seam Tracking and Detection Robot Based on Artificial Intelligence Technology[J]. Sensors, 2023, 23(15): 6725. [14] Du Y, Liu M, Wang J, et al. A wall climbing robot based on machine vision for automatic welding seam inspection[J]. Ocean Engineering, 2024, 310: 118825. [15] Li M, Huang J, Xue L, Zhang R. A guidance system for robotic welding based on an improved YOLOv5 algorithm with a RealSense depth camera[J]. Scientific Reports, 2023, 13: 21299. [16] Zhu S, Zhang L, Zhao J, Hu X, Luo X. Research on steel structure weld seam recognition algorithm based on improved YOLOv5[J]. Scientific Reports, 2025, 15: 34276. [17] Wang S, Li Z, Chen G, Yue Y. Weld seam object detection system based on the fusion of 2D images and 3D point clouds using interpretable neural networks[J]. Scientific Reports, 2024, 14: 21137. [18] Pham D A, Bui D Q, Le T D, Tran D H, Nguyen T H. Automatic welding seam tracking and real-world coordinates identification with machine learning method[J]. Results in Engineering, 2024, 23: 102565. [19] 邢正, 李斌, 梁志达, 等. 球罐爬壁机器人的焊缝双向跟踪方法研究[J]. 机械设计与制造, 2026(2): 342-346. [20] 邓弘文, 石海峡, 朱利凯, 等. 优化模糊控制算法的四轮差速机器人运动控制[J]. 机械研究与应用, 2025, 38(2): 7-13. [21] 张树晓. PSO优化模糊PID的光伏板清洁机器人防偏摆运动控制研究[J]. 机械设计与制造工程, 2025, 54(2): 47-51. [22] 邱东. 融合Frenet坐标系的室内寻迹机器人运动控制优化研究[D]. 南京邮电大学, 2023.\\

\chapter*{第二章 视觉感知与目标位置表征方法}

本章围绕船体表面焊缝跟踪任务中的视觉感知问题展开。面向爬壁机器人焊缝作业背景，系统首先需要从图像中获得焊缝目标的位置；其次，需要将该位置由图像检测结果转化为控制系统可以使用的几何量；最后，需要根据目标位置与期望参考位置之间的关系形成偏差信息。针对上述过程，本文采用基于目标检测的焊缝区域识别方法，以检测框中心作为焊缝目标位置表征，并以图像中心作为期望参考建立横向偏差，为第三章运动控制方法提供输入。\\

\section{2.1 焊缝跟踪任务中的视觉感知需求}

焊缝跟踪任务本质上要求机器人根据目标在视野中的位置变化实时调整运动方向。对于搭载相机的爬壁机器人或移动执行平台而言，外部环境首先以二维图像的形式进入系统。图像中的焊缝目标可能表现为一段线状区域、局部纹理区域或具有一定宽度的连续目标区域。控制系统并不能直接使用原始图像像素生成运动指令，而需要一个能够表达目标相对位置的中间量。因此，视觉感知环节需要完成从“图像信息”到“目标位置量”的转换。\\

在传统图像处理方法中，常通过颜色阈值、边缘提取、形态学处理或扫描线分析获得目标区域。这类方法结构较简单，但容易受到背景颜色、光照变化和图像噪声影响。本文系统保留了基于 HSV 阈值的传统视觉路径作为备用方式，但主运行链路采用 YOLOv5 目标检测方法输出焊缝目标检测框，并进一步提取检测框中心作为控制依据。因此，本文视觉感知部分不把 HSV 阈值分割写成主方法，而将其作为传统方法背景。\\

焊缝跟踪对视觉输出有两个要求。第一，视觉输出必须能够反映目标在图像中的横向位置，因为机器人转向控制主要依赖目标相对于图像中心的左右偏移。第二，视觉输出应当尽量简洁，避免将完整图像或复杂检测结果直接传递给控制逻辑。基于这两个要求，本文采用检测框中心作为焊缝目标的位置表征。该方法不需要提取完整焊缝中心线，也不需要进行三维重建，而是通过检测框的左右边界计算目标横向中心，从而建立视觉感知与运动控制之间的接口。\\

需要说明的是，本文视觉方法面向焊缝跟踪控制所需的横向位置表征展开，不讨论焊缝语义分割、骨架提取、中心线拟合、焊缝宽度测量或三维重建等超出当前系统范围的视觉任务。因此，本章只围绕目标检测框、目标中心、图像中心和横向偏差展开。\\

\section{2.2 图像输入与视觉检测流程}

视觉检测流程的第一步是获得稳定的图像输入。本文系统可以接收真实相机图像，也可以使用本地图片或视频帧作为测试输入。不同图像来源在进入视觉检测流程后被统一处理，这使得检测方法不依赖具体采集方式。对于论文方法表述而言，可以将其概括为统一图像输入接口：无论图像来自相机还是离线测试素材，后续处理流程均围绕同一类图像数据展开。\\

图像进入系统后，需要完成格式转换和模型预处理。系统将输入图像转换为 OpenCV 图像格式，并按照 YOLOv5 推理流程进行尺寸调整、颜色通道转换、维度变换、张量化、浮点化和归一化处理。其中，尺寸调整采用 YOLOv5 工具链中的 letterbox 方法，其作用是在适配模型输入尺寸的同时尽量保持原始图像比例。经过预处理后，图像被送入目标检测模型进行推理，模型输出候选检测框。\\

检测模型输出后，系统还需要对候选框进行筛选。本文采用置信度阈值和非极大值抑制对候选框进行过滤，并将检测框坐标恢复到原始图像坐标系中。坐标恢复是视觉感知与控制输入连接的重要环节，因为后续图像中心和目标中心都需要在同一坐标系下计算。如果检测框仍停留在模型输入尺寸坐标系中，目标中心与原始图像中心之间的偏差就不具备统一尺度，控制输出也会受到影响。\\

因此，本文将视觉检测流程概括为：图像输入、图像预处理、目标检测推理、候选框筛选、坐标恢复和目标位置输出。该流程的输入是相机或测试素材提供的图像，输出是焊缝目标检测框及其有效性状态。它为后续目标中心提取和偏差构建奠定了基础。\\

\section{2.3 基于目标检测的焊缝区域识别方法}

本文采用基于目标检测的焊缝区域识别方法。目标检测方法的输出通常由目标框位置、置信度和类别信息组成。对于焊缝跟踪任务而言，最关键的是目标框在图像中的几何位置。目标框能够描述焊缝目标所在的图像区域，进而可以通过几何计算得到目标中心位置。\\

系统检测过程可表示为输入图像 $I_t$ 到检测结果 $D_t$ 的映射：\\

\[
D_t = F_{\text{det}}(I_t)
\]
其中，$I_t$ 表示 $t$ 时刻输入图像，$F_{\text{det}}(\cdot)$ 表示目标检测推理过程，$D_t$ 表示模型输出并经过筛选后的检测结果。该公式是对检测流程的理论化表达，用于说明视觉感知环节的输入和输出关系。本文中的 $F_{\text{det}}(\cdot)$ 表示已加载权重后的 YOLOv5 推理过程，不表示本文重新设计或训练了新的网络结构。\\

单个检测框可表示为：\\

\[
B_t = (x_{\min},y_{\min},x_{\max},y_{\max},c_t)
\]
其中，$B_t$ 表示 $t$ 时刻选取的目标检测框，$x_{\min}$、$y_{\min}$ 表示检测框左上角坐标，$x_{\max}$、$y_{\max}$ 表示检测框右下角坐标，$c_t$ 表示检测置信度。图像坐标系以图像左上角为原点，$x$ 轴向右为正，$y$ 轴向下为正。该表达式用于将目标检测输出抽象为几何边界框与置信度。系统选择最高置信度检测框作为当前帧控制目标，因此 $c_t$ 在目标选择中具有作用；但后续控制主要使用检测框横向中心，并未使用置信度参与速度计算。\\

采用目标检测框作为焊缝区域表征，符合当前系统的实现边界。检测框能够提供目标区域的空间范围，却不要求系统获得每一个焊缝像素点的位置。对于以横向偏差为控制输入的跟踪系统而言，检测框中心已经可以反映目标相对图像中心的位置关系。因此，本文不引入复杂的焊缝轮廓提取和中心线拟合，而将检测框作为视觉识别与控制之间的中间表示。\\

需要注意的是，目前尚未整理形成权重对应的训练数据集、类别名称、训练轮次、网络规模、损失曲线和 mAP 指标。因此，本文不写“模型训练取得某精度”“检测准确率达到某数值”或“相较其他模型有明显提升”等缺少实验支撑的结论。若后续补充训练日志和验证结果，可在本节之后增加模型训练与评价小节。\\

\section{2.4 焊缝目标中心位置表征方法}

目标检测框给出了焊缝目标区域的边界，但运动控制并不需要直接使用完整边界框。对于当前系统而言，控制部分主要关心目标在横向方向上相对图像中心的位置，因此需要将检测框进一步压缩为目标中心位置。\\

设检测框左、右边界横坐标分别为 $x_{\min}$、$x_{\max}$，上、下边界纵坐标分别为 $y_{\min}$、$y_{\max}$，则检测框中心点可表示为：\\

\[
x_c = \frac{x_{\min}+x_{\max}}{2}
\]
\[
y_c = \frac{y_{\min}+y_{\max}}{2}
\]
其中，$x_c$ 表示焊缝目标在图像坐标系中的横向中心位置，$y_c$ 表示焊缝目标在图像坐标系中的纵向中心位置。该公式用于将目标检测框转换为中心点坐标。本文控制逻辑实际使用的是横向中心 $x_c$，因为机器人转向主要由目标相对于图像中心的左右偏移决定。纵向中心 $y_c$ 可以作为检测框几何中心的理论表达，但不参与当前速度控制量生成。\\

目标中心提取的作用可以从感知和控制两个角度理解。从感知角度看，它将检测模型输出的矩形区域转化为一个代表性点，使目标位置表达更简洁。从控制角度看，它将视觉结果转换为后续偏差计算所需的坐标量，使控制器无需处理完整图像和完整检测框。由此，视觉模块输出从复杂图像数据变成了目标横向中心、图像宽度和有效性标志。\\

为了表达目标中心输出，本文将 $t$ 时刻视觉位置表征写为：\\

\[
p_t = (x_c,W,s_t)
\]
其中，$p_t$ 表示视觉模块输出给控制模块的位置表征，$W$ 表示图像宽度，$s_t$ 表示检测结果有效性标志。当检测有效时，$s_t=1$；当检测无效时，$s_t=0$。在 ROS 实现中，该位置表征可通过包含三个数值分量的消息传递，其中三个分量分别对应目标中心、图像全宽和有效标志。该表达使视觉检测结果能够以简洁形式进入后续控制环节。\\

有效性标志可进一步写成分段形式：\\

\[
s_t =
\begin{cases}
1, & \text{检测结果有效} \\
0, & \text{检测结果无效}
\end{cases}
\]
该公式说明视觉输出不仅包含目标位置，也包含检测是否可信的状态信息。对于焊缝跟踪任务而言，这一点十分重要。若系统在目标丢失时继续使用上一帧目标位置，机器人可能产生错误运动；而有效性标志可以使控制部分在检测失败时及时停止。\\

根据当前系统逻辑，视觉输出可进一步抽象为：\\

\[
p_t =
\begin{cases}
(x_c,W,1), & s_t=1 \\
(-1,W,0), & s_t=0
\end{cases}
\]
其中，$(-1,W,0)$ 表示无效检测情况下的位置输出形式。该表达式连接了视觉检测状态和控制安全策略，为第三章安全停止逻辑提供输入。\\

\section{2.5 图像偏差构建与控制输入转换}

完成目标中心位置表征后，需要进一步建立目标中心与期望参考位置之间的关系。对于基于前视图像的焊缝跟踪任务，可以将图像中心线作为期望参考。若焊缝目标中心位于图像中心附近，说明机器人当前运动方向与目标位置较为一致；若目标中心偏离图像中心，则说明机器人需要调整方向。\\

设图像宽度为 $W$，高度为 $H$，则图像中心横坐标和纵坐标可表示为：\\

\[
x_0 = \frac{W}{2}
\]
\[
y_0 = \frac{H}{2}
\]
其中，$x_0$ 表示图像中心横坐标，$y_0$ 表示图像中心纵坐标。当前系统主要使用 $x_0$ 作为横向参考位置。图像中心参考的作用是建立期望跟踪位置：焊缝目标中心越接近 $x_0$，机器人横向偏差越小；焊缝目标中心越远离 $x_0$，机器人需要越明显地调整运动方向。\\

焊缝目标中心与图像中心之间的像素横向偏差可表示为：\\

\[
e_p = x_c - x_0
\]
其中，$e_p$ 为像素坐标下的横向偏差。$e_p$ 的正负号表示焊缝目标相对图像中心的左右偏移方向，绝对值表示偏移程度。该公式连接了视觉目标位置表征与运动控制输入，是视觉结果进入控制环节的关键步骤。需要注意的是，$e_p$ 的符号方向与机器人最终转向方向之间还受到相机安装方向、图像坐标系方向和控制坐标约定影响，因此实际控制方向应以运行现象为准。\\

为了降低图像分辨率变化对控制量尺度的影响，可以将像素偏差归一化：\\

\[
e_n = \frac{e_p}{W/2}
\]
将 $e_p=x_c-W/2$ 代入上式，可得：\\

\[
e_n = \frac{x_c-W/2}{W/2}
\]
其中，$e_n$ 表示归一化横向偏差。该量将偏差转换为相对于图像半宽的比例关系，使偏差表达更适合用于速度控制。本文将 $e_n$ 作为视觉偏差进入控制环节前的理论抽象，用于说明不同图像宽度下偏差尺度统一的必要性。\\

视觉输出到控制输入的转换关系可写为：\\

\[
e_t = g(p_t)
\]
其中，$g(\cdot)$ 表示由视觉位置表征到控制偏差的转换函数。若只从几何位置关系描述目标相对图像中心的偏移，可写为：\\

\[
e_t = \frac{x_c-W/2}{W/2},\quad s_t=1
\]
上述 $e_t$ 与 $e_n$ 方向一致，用于说明焊缝目标相对图像中心的几何偏移。第三章控制器内部实际采用相反方向的控制误差：\\

\[
e = \frac{W/2-x_c}{W/2} = -e_n
\]
因此，$e_n$ 用于描述图像几何偏差，$e$ 用于生成角速度控制量，二者只差一个符号方向。\\

在检测无效时，控制系统不再根据该帧位置计算跟踪速度，而应进入停止逻辑：\\

\[
e_t \ \text{无效},\quad s_t=0
\]
由此可见，视觉感知模块不仅提供目标位置，还提供控制是否继续执行的判断条件。该设计使系统能够在目标丢失、检测失败或输入异常时避免继续按照旧偏差运动，从而提高运行安全性。第二章建立的 $x_c$、$x_0$、$e_p$ 和 $e_n$ 将在第三章中进一步用于角速度和线速度控制量生成。\\

\section{2.6 本章小结}

本章围绕焊缝跟踪任务中的视觉感知与目标位置表征方法进行了阐述。系统首先通过统一图像输入获取相机或测试视频图像，然后采用 YOLOv5 目标检测方法识别焊缝区域，并通过检测框左右边界均值计算目标横向中心。在此基础上，本文以图像中心作为期望参考，构建像素横向偏差和归一化偏差，为后续运动控制提供输入。\\

本章中的检测框中心公式 $x_c=(x_{\min}+x_{\max})/2$ 和图像中心公式 $x_0=W/2$ 对应本文视觉位置表征方法；归一化偏差公式 $e_n=(x_c-W/2)/(W/2)$ 是对比例偏差计算的论文式抽象；有效性标志 $s_t$ 反映了检测结果能否作为控制输入。目前实验材料不能支撑训练数据集、检测精度、mAP、语义分割、中心线拟合或三维重建等内容，相关内容需要后续补充证据后才能写入正式论文。\\

\section{本章依据说明}

本章依据系统实现资料中的“基于 ROS 图像话题的视觉输入方法”“基于 YOLOv5 的目标检测推理方法”“检测框中心位置表征方法”“目标丢失与输入超时安全停止方法”和“训练、评估和可视化工具的存在边界”。具体实现依据包括图像输入、YOLOv5 推理、检测框中心计算、中心点有效标志和结果图可视化，但正文未将具体路径作为主体叙述。\\

\chapter*{第三章 运动控制与焊缝跟踪方法}

第二章建立了焊缝目标检测框、目标中心、图像中心和横向偏差之间的关系。本章在此基础上讨论运动控制问题，即如何将视觉偏差信息转化为机器人速度控制量。当前系统面向麦克纳姆轮全向移动底盘的速度接口，采用以偏差抑制、运动平稳和安全执行为目标的视觉闭环跟踪控制方法，控制目标是使焊缝目标尽量保持在图像中心附近，并在目标丢失或输入超时时停止运动。\\

为了体现焊缝跟踪过程中的运动优化，本文不把控制部分仅表述为简单的速度输出，而是围绕视觉闭环建立分层运动优化策略。该策略首先利用全向底盘速度接口建立车体速度层运动模型，然后将检测框中心偏差归一化为控制误差，在此基础上引入小增益积分修正、积分分离、积分限幅以及线速度-角速度耦合约束，并结合目标丢失停止机制形成完整的局部跟踪控制链路。由此，运动优化不再停留在单一转向控制上，而是落实为误差建模、角速度生成、线速度分配和失效保护共同参与的控制策略设计。\\

\section{3.1 焊缝跟踪控制任务分析}

焊缝跟踪控制的任务是根据视觉系统输出的目标位置信息生成机器人运动指令，使机器人在运动过程中不断修正自身运动状态。对于本文系统而言，视觉模块输出的是目标中心 $x_c$、图像宽度 $W$ 和检测有效标志 $s_t$。控制模块并不直接处理原始图像，也不直接处理完整检测框，而是根据目标中心与图像中心之间的关系计算运动控制量。\\

控制任务可以分为四个层次。第一是偏差判断，即判断焊缝目标位于图像中心左侧还是右侧，以及偏移程度有多大。第二是控制误差构建，即将像素尺度下的目标偏移转换为归一化误差，减少图像宽度变化对控制量的影响。第三是速度生成与约束，即根据偏差方向生成角速度，根据角速度和偏差状态约束前进速度。第四是异常处理，即当检测结果无效或输入长时间中断时，停止机器人运动并清除积分残余。四个层次共同构成了本文运动优化控制策略的基础。\\

在焊缝跟踪过程中，若目标中心接近图像中心，说明机器人运动方向与焊缝目标方向较为一致，此时可保持前进速度；若目标中心明显偏离图像中心，说明机器人需要转向修正；若目标无法检测，则当前视觉信息不足以支撑继续运动，系统应进入停止状态。该控制思想结构相对简洁，适合在缺少复杂状态反馈的条件下完成基础焊缝跟踪。\\

\section{3.2 麦克纳姆轮全向底盘速度接口与运动学模型}

麦克纳姆轮底盘通过斜向滚子与四轮组合，使机器人具备平面内前后移动、横向移动和原地旋转的运动能力。与普通两轮驱动底盘相比，全向移动底盘在理论上具有更高的运动自由度，能够通过车体坐标系下的 $v_x$、$v_y$ 和 $\omega$ 描述平面三自由度速度。其中，$v_x$ 表示车体前向速度，$v_y$ 表示车体横向速度，$\omega$ 表示绕竖直轴的偏航角速度。\\

设机器人在全局平面坐标系中的位姿为：\\

\[
\mathbf{q}=[x,\ y,\ \theta]^T
\]
其中，$x$ 和 $y$ 表示机器人在平面坐标系中的位置，$\theta$ 表示车体航向角。设车体坐标系下速度输入为：\\

\[
\mathbf{u}=[v_x,\ v_y,\ \omega]^T
\]
则全向底盘在车体速度层的平面运动学关系可表示为：\\

\[
\begin{bmatrix}
\dot{x}\\
\dot{y}\\
\dot{\theta}
\end{bmatrix}
=
\begin{bmatrix}
\cos\theta & -\sin\theta & 0\\
\sin\theta & \cos\theta & 0\\
0 & 0 & 1
\end{bmatrix}
\begin{bmatrix}
v_x\\
v_y\\
\omega
\end{bmatrix}
\]
展开可得：\\

\[
\dot{x}=v_x\cos\theta-v_y\sin\theta
\]
\[
\dot{y}=v_x\sin\theta+v_y\cos\theta
\]
\[
\dot{\theta}=\omega
\]
该模型说明了全向底盘车体速度指令与平面位姿变化之间的关系。当 $v_y$ 不为零时，机器人可以在不改变航向的情况下产生横向位移；当 $\omega$ 不为零时，机器人发生偏航旋转。当前底盘速度接口能够接收 $v_x$、$v_y$ 和 $\omega$ 三个速度分量，因此可用上述模型解释全向移动平台的速度输入含义。\\

\subsection{3.2.1 底盘执行层速度分配与电机闭环实现}

在车体速度层之外，底盘执行层还需要将 $v_x$、$v_y$ 和 $\omega$ 转换为四个电机的目标转速。根据 NanoOmni 控制板资料，底盘通信协议中速度控制指令的功能码为 $0x01$，数据区由 X 轴线速度、Y 轴线速度和 Z 轴角速度三个分量组成，三个量均按 1000 倍放大后以有符号 16 位整数发送。上位机侧的底盘接口读取标准速度指令中的前向速度、横向速度和偏航角速度，并按该协议发送到底盘控制板；控制板接收到速度帧后再将三个速度分量恢复为实际速度量。\\

设轮距为 $T$，轴距为 $B$，轮径为 $D$，四个电机目标转速为 $n_1,n_2,n_3,n_4$，单位为 r/min。结合控制板中麦克纳姆轮速度分配关系，车体速度到四轮目标转速的转换可写为：\\

\[
n_1=\frac{60\left[v_x+v_y-\omega(T+B)/2\right]}{\pi D}
\]
\[
n_2=\frac{60\left[v_x-v_y+\omega(T+B)/2\right]}{\pi D}
\]
\[
n_3=\frac{60\left[v_x-v_y-\omega(T+B)/2\right]}{\pi D}
\]
\[
n_4=\frac{60\left[v_x+v_y+\omega(T+B)/2\right]}{\pi D}
\]
其中，系数 60 用于将 m/s 和 rad/s 对应的速度量转换到每分钟尺度，再除以车轮周长 $\pi D$ 得到电机目标转速。该公式的符号方向与控制板中的电机编号、车体坐标系和麦克纳姆轮安装方向对应。根据 NanoOmni 控制板资料，本文建模暂按 $T=0.168$ m，$B=0.145$ m，$D=0.064$ m 取值，而不将其表述为实测尺寸，则有：\\

\[
\frac{60}{\pi D}\approx 298.42
\]
\[
\frac{T+B}{2}=0.1565
\]
因此，四轮目标转速也可近似写为：\\

\[
n_1\approx 298.42(v_x+v_y-0.1565\omega)
\]
\[
n_2\approx 298.42(v_x-v_y+0.1565\omega)
\]
\[
n_3\approx 298.42(v_x-v_y-0.1565\omega)
\]
\[
n_4\approx 298.42(v_x+v_y+0.1565\omega)
\]
上述关系说明，底盘执行层不仅接收车体速度指令，而且会在控制板侧完成四轮目标转速分配。控制板资料还给出了减速比、编码器和电机控制参数：减速比为 21.3，编码器线数为 11，电机控制周期为 25 ms，电机速度闭环参数为 $K_P=0.2$、$K_I=0.2$、$K_D=0.25$。在执行过程中，控制板通过四路编码器读取电机转动增量并计算当前转速，再根据目标转速与实际转速误差进行电机 PID 计算，最后输出 PWM 驱动电机。由此可见，本文上层焊缝跟踪控制输出的是车体速度指令，而四轮分配、编码器测速和电机速度闭环属于底盘执行层完成的内容。\\

需要说明的是，上述底盘参数来自本地 NanoOmni 控制板资料及其配套程序。本文在运动学建模和执行链路说明中先按该组资料参数取值，用于支撑底盘执行层的公式推导与参数化说明；但该组参数仍需与实际车辆固件版本、拨码配置和实物尺寸进一步确认一致性，因此不写成实测几何尺寸。\\

当前焊缝跟踪控制实际采用前向速度与偏航角速度进行局部跟踪，横向速度通道保留但未用于纠偏。因此，本文控制输入可写为：\\

\[
\mathbf{u}_k=[v_k,\ 0,\ \omega_k]^T
\]
其中，$v_k$ 表示第 $k$ 次控制输出中的前向速度指令值，$\omega_k$ 表示偏航角速度指令值。代入全向底盘车体速度层模型后，可得到当前实现对应的简化关系：\\

\[
\dot{x}=v_k\cos\theta
\]
\[
\dot{y}=v_k\sin\theta
\]
\[
\dot{\theta}=\omega_k
\]
该简化模型并不否定麦克纳姆轮底盘的全向运动能力，而是说明当前焊缝跟踪算法在上层控制中选择了更保守的前向速度加偏航修正方式。这样既能够保持原有控制链路的稳定性，又能通过底盘三轴速度接口保留后续扩展横向纠偏的可能。\\

同时，将 $v_y=0$ 代入底盘执行层轮速分配公式，可得当前焊缝跟踪策略下的四轮目标转速形式：\\

\[
n_1=\frac{60\left[v_k-\omega_k(T+B)/2\right]}{\pi D}
\]
\[
n_2=\frac{60\left[v_k+\omega_k(T+B)/2\right]}{\pi D}
\]
\[
n_3=\frac{60\left[v_k-\omega_k(T+B)/2\right]}{\pi D}
\]
\[
n_4=\frac{60\left[v_k+\omega_k(T+B)/2\right]}{\pi D}
\]
该式表明，在当前焊缝跟踪控制中，同侧或对角轮速差主要由偏航角速度项决定，机器人通过前向运动叠加转向修正来减小图像中心偏差；横向速度 $v_y$ 虽然是底盘能力的一部分，但没有被当前上层偏差控制策略用于横向平移纠偏。\\

\section{3.3 基于视觉偏差的跟踪误差构建}

运动控制的输入来自视觉偏差。根据第二章定义，检测框横向中心为 $x_c$，图像中心横坐标为 $x_0=W/2$，像素横向偏差为：\\

\[
e_p = x_c - x_0
\]
其中，$e_p$ 表示焊缝目标中心相对于图像中心的像素偏移量。若 $e_p$ 接近零，说明目标中心接近图像中心；若 $e_p$ 偏离零，说明目标位于图像中心左侧或右侧。该偏差是视觉识别结果进入控制环节的直接依据。\\

由于图像分辨率可能变化，仅使用像素偏差会使控制量尺度受到图像宽度影响。为便于建立统一控制关系，本文使用归一化偏差：\\

\[
e_n = \frac{e_p}{W/2}
\]
即：\\

\[
e_n = \frac{x_c-W/2}{W/2}
\]
其中，$e_n$ 表示相对于图像半宽的横向偏差比例。该量的绝对值越大，说明焊缝目标偏离图像中心越明显；其符号表示偏移方向。归一化处理是本文运动优化策略中的第一层，其作用是将不同图像宽度下的像素偏差统一到相同量纲的比例误差中，使后续控制参数不直接依赖像素宽度。\\

当前控制函数内部采用的误差方向为图像中心减目标中心，即：\\

\[
e_k = \frac{x_0-x_c}{x_0}
\]
结合 $x_0=W/2$ 和 $e_n=(x_c-x_0)/x_0$ 可得：\\

\[
e_k = -e_n
\]
因此，$e_n$ 用于描述目标相对图像中心的几何偏移，$e_k$ 用于表示控制器内部的修正误差。采用相反符号的原因是控制输出需要使目标中心向图像中心靠近，即角速度方向应与几何偏差方向构成负反馈关系。符号方向与图像坐标系、相机安装方向、底盘坐标系和控制约定有关，正式实机解释应以运行现象为准。\\

为了在控制中区分目标是否接近中心，可设定允许偏差阈值 $\varepsilon$。当满足：\\

\[
|e_k| < \varepsilon
\]
时，可认为焊缝目标位于图像中心附近，机器人可以保持前进；当 $|e_k|\geq\varepsilon$ 时，控制系统需要根据偏差方向生成转向动作。本文取 $\varepsilon=0.05$，表示目标中心偏差小于图像半宽的 5\% 时进入中心死区。以仿真相机常用图像宽度 $W=640$ px 为例，$W/2=320$ px，因此该死区约对应图像中心左右 $\pm 16$ px。\\

\section{3.4 面向视觉闭环的运动优化控制策略}

视觉焊缝跟踪中的运动优化并不只等同于提高速度，而是需要同时考虑偏差减小、转向平稳、速度受限和异常停止。为说明控制策略的设计目标，可将单步运动优化目标抽象为：\\

\[
J_k=\lambda_e |e_k|+\lambda_\omega |\omega_k|+\lambda_v |v_k-v_{k-1}|+\lambda_s(1-s_k)
\]
其中，$|e_k|$ 表示跟踪偏差，$|\omega_k|$ 反映转向强度，$|v_k-v_{k-1}|$ 反映速度变化平滑性，$1-s_k$ 表示检测无效带来的安全约束项，$\lambda_e$、$\lambda_\omega$、$\lambda_v$ 和 $\lambda_s$ 为权重系数。该表达式用于说明本文运动优化所关注的目标，而不是表示系统采用在线数值优化求解器。本文采用参数化控制律对上述目标进行工程化近似，即通过归一化误差、PI 修正、速度耦合和停止机制实现对运动过程的约束。\\

\subsection{3.4.1 小增益 PI 偏差修正}

角速度用于修正目标横向偏差。若只采用比例控制，当检测框中心长期存在小偏差时，系统可能出现残余误差。为减小纯比例控制可能存在的残余偏差，本文在原有比例项基础上引入小增益积分项。定义控制误差为 $e_k=(x_0-x_c)/x_0$，设第 $k$ 次控制计算的时间间隔为 $\Delta t_k$，并对异常时间间隔进行保护：\\

\[
\Delta t_k =
\begin{cases}
t_k-t_{k-1}, & 0<t_k-t_{k-1}\leq 1.0\\
0.1, & t_k-t_{k-1}\leq 0\ \text{或}\ t_k-t_{k-1}>1.0
\end{cases}
\]
积分状态 $I_e$ 的更新规则为：\\

\[
I_e[k]=
\begin{cases}
\operatorname{clip}\left(I_e[k-1]+e_k\Delta t_k,\ -I_{\max},\ I_{\max}\right), & |e_k|<e_{\text{sep}}\\[4pt]
\operatorname{clip}\left(I_e[k-1],\ -I_{\max},\ I_{\max}\right), & |e_k|\geq e_{\text{sep}}
\end{cases}
\]
其中，$I_{\max}=0.3$ 为积分限幅，$e_{\text{sep}}=0.30$ 为积分分离阈值。二者取值相近但含义不同：前者限制积分状态的最大幅值，后者决定是否继续累积积分。积分分离的作用是在偏差过大时避免积分项快速累积，在偏差进入较小范围后再利用积分项修正残余误差。\\

角速度控制律为：\\

\[
\omega_k =
\begin{cases}
0, & |e_k|<\varepsilon\\[4pt]
K_p e_k + K_i I_e[k], & |e_k|\geq\varepsilon
\end{cases}
\]
其中，$\omega_k$ 表示第 $k$ 次控制输出的角速度，$K_p$ 为比例增益，$K_i$ 为积分增益。本文取 $K_p=0.5$、$K_i=0.02$。$K_i$ 取值较小，用于在尽量保持原有比例控制动态特性的基础上引入弱积分作用；微分项未引入，因此本文不将该控制器写成完整 PID 控制器。死区内角速度为零，积分状态不因死区而继续累加；检测无效或输入超时时，积分状态会被清零。\\

\subsection{3.4.2 线速度-角速度耦合约束}

线速度用于保持机器人沿焊缝目标方向前进。若机器人在偏差较大时仍以较高速度前进，目标可能进一步偏离视野；若机器人在偏差较小时速度过低，则跟踪效率不足。因此，本文将前进速度与转向强度进行耦合约束。\\

当目标接近图像中心时，即 $|e_k|<\varepsilon$ 且 $\varepsilon=0.05$，系统输出直行速度指令值，角速度保持为零：\\

\[
(\omega_k,v_k)=(0,\ v_0),\quad v_0=0.2,\quad |e_k|<0.05
\]
其中，$v_0=0.2$ 表示前向速度指令值。ROS 标准速度指令的名义单位为 m/s，但目前尚未提供实机标定关系，因此本文将其表述为速度指令值，而不写成已经实测确认的物理速度。\\

当目标偏差超过死区范围时，线速度与角速度耦合以兼顾稳定性和跟踪效率：\\

\[
v_k =
\begin{cases}
v_0, & |e_k|<\varepsilon \\[6pt]
\max\!\Big(v_{\min},\ v_0 - \alpha |\omega_k|\Big), & |e_k|\geq\varepsilon
\end{cases}
\]
其中，$v_0=0.2$ 为直行速度指令值，$v_{\min}=0.1$ 为最低速度指令值，$\alpha=0.5$ 为线速度-角速度耦合系数。系统先判断 $|\omega_k|<0.2$，若满足则令 $v_k=0.2-|\omega_k|/2$，否则令 $v_k=0.1$。由于 $|\omega_k|=0.2$ 时 $0.2-|\omega_k|/2=0.1$，上述逻辑也可写成 $v_k=\max(0.1,0.2-|\omega_k|/2)$。这一设计使机器人在转向较大时自动降低前进速度，减小目标丢失风险。\\

\subsection{3.4.3 控制参数与优化含义}

本文运动控制参数可整理如下。$K_p=0.5$ 为比例增益，用于根据当前偏差生成主要角速度修正量；$K_i=0.02$ 为积分增益，用于对残余偏差进行小幅积分修正；$\varepsilon=0.05$ 为中心死区阈值，用于判断目标是否接近图像中心；$I_{\max}=0.3$ 为积分限幅，用于限制积分项幅值并降低积分过大风险；$e_{\text{sep}}=0.30$ 为积分分离阈值，表示系统只在较小偏差范围内累积积分；$v_0=0.2$ 为直行速度指令值，对应目标接近中心时的前进速度；$v_{\min}=0.1$ 为最低速度指令值，用于在转向较大时保留较低前进速度；$\alpha=0.5$ 为速度耦合系数，表示角速度增大时线速度降低的比例；$T_{\text{out}}=0.5$ s 为中心点超时阈值，用于在输入中断时触发停止。\\
由此可见，本文的运动优化不是孤立的参数罗列，而是围绕视觉闭环跟踪目标构建的控制策略组合。归一化误差解决像素尺度问题，小增益 PI 修正解决残余偏差问题，积分分离和限幅降低积分过大风险，线速度-角速度耦合改善转向时的运动平稳性，安全停止机制限制目标丢失后的误运动风险。这些设计共同服务于焊缝目标保持和局部跟踪稳定性。\\

\section{3.5 感知—控制闭环机制}

焊缝跟踪系统不是单向检测过程，而是感知与控制之间的循环关系。输入图像经过检测获得目标位置，目标位置经过偏差计算生成控制量，控制量作用于机器人后会改变下一时刻图像中的目标位置。该闭环关系可表示为：\\

\[
I_t \rightarrow D_t \rightarrow p_t \rightarrow e_t \rightarrow u_t \rightarrow q_{t+1}
\]
其中，$I_t$ 表示输入图像，$D_t$ 表示检测结果，$p_t$ 表示目标位置表征，$e_t$ 表示偏差信息，$u_t=[v_t,\ 0,\ \omega_t]^T$ 表示当前焊缝跟踪控制输出，$q_{t+1}$ 表示机器人下一时刻运动状态。该公式说明视觉感知与运动控制之间的信息传递关系。\\

在该闭环中，视觉感知环节负责回答“焊缝目标在图像中的什么位置”；位置表征环节负责将检测框转换为目标中心；偏差构建环节负责判断目标相对期望中心的偏移；运动优化控制环节负责在偏差减小、转向平稳和速度受限之间进行折中；执行环节负责使机器人或仿真模型产生运动响应。通过该链路，机器人能够根据视觉结果不断调整运动。\\

本文闭环属于基于视觉输入的局部控制闭环，而不是基于全局地图、里程计反馈或多传感器融合的导航闭环。系统中虽然存在底盘里程计发布和其他传感器相关功能，但焊缝跟踪控制逻辑没有消费这些反馈。因此，本文不将该闭环写成完整定位导航闭环，而写成视觉偏差驱动的局部跟踪闭环。\\

\section{3.6 异常处理与安全停止逻辑}

在实际运行中，视觉检测可能出现无目标、目标被遮挡、检测置信度不足或图像输入中断等情况。如果控制系统在这些情况下继续使用上一帧偏差，机器人可能沿错误方向运动。因此，系统需要异常处理机制，使控制输出在视觉结果无效时进入安全状态。\\

当前系统通过检测有效标志和输入超时判断实现停止逻辑。当检测有效时，控制系统根据偏差生成速度；当检测无效时，控制系统输出零速度。该关系可写为：\\

\[
\mathbf{u}_t =
\begin{cases}
[v_t,\ 0,\ \omega_t]^T, & s_t=1 \\
[0,\ 0,\ 0]^T, & s_t=0
\end{cases}
\]
其中，$\mathbf{u}_t$ 表示控制输出，$s_t$ 表示视觉结果有效性标志。当 $s_t=1$ 时，系统使用视觉偏差生成速度控制量；当 $s_t=0$ 时，系统输出零速度。此外，检测无效时积分累积量 $I_e$ 被清零，防止恢复跟踪后产生与停止前偏差方向无关的积分残余：\\

\[
I_e \leftarrow 0,\quad s_t=0
\]
此外，系统还设置了输入超时处理。若长时间未收到有效目标中心，控制部分同样输出零速度。可将超时条件表示为：\\

\[
\Delta t = t - t_{\text{last}}
\]
\[
\mathbf{u}_t=[0,\ 0,\ 0]^T,\quad \Delta t > T_{\text{out}}
\]
其中，$t_{\text{last}}$ 表示最近一次收到有效视觉输入的时刻，$T_{\text{out}}$ 表示允许的超时时间。本文取 $T_{\text{out}}=0.5$ s，并在超时停止时同步清零积分状态：\\

\[
I_e \leftarrow 0,\quad \Delta t > T_{\text{out}}
\]
该公式用于描述输入中断情况下的停止逻辑。安全停止逻辑的意义在于限制系统在感知不可靠时继续运动。它不能替代硬件急停或工业安全认证，也不能说明系统已经完成实机安全测试；后续仍需通过实机运行视频和速度输出记录验证停止效果。\\

\section{3.7 本章小结}

本章围绕运动控制与焊缝跟踪方法进行了阐述。首先，结合麦克纳姆轮全向移动底盘的速度接口，建立了车体速度层平面运动学模型，并进一步说明底盘执行层如何将三轴速度指令分配为四个电机目标转速，以及如何通过编码器测速和电机 PID 闭环完成执行。其次，本文明确当前焊缝跟踪控制实际采用 $\mathbf{u}_k=[v_k,\ 0,\ \omega_k]^T$ 的简化输入，并基于检测框中心与图像中心之间的横向偏差构建归一化控制误差。最后，围绕视觉闭环提出由小增益 PI 修正、积分分离、积分限幅、线速度-角速度耦合约束和安全停止机制组成的运动优化控制策略。\\

本章的运动优化特点是在保持原有控制结构和接口稳定的基础上，对误差构建、角速度生成、速度约束和异常处理进行参数化改进。换言之，本文在运动控制部分的主要工作量体现在将视觉偏差闭环细化为可分析、可调参与可验证的策略组合，而不是只给出单一比例转向关系。当前可确认的是系统具备全向底盘三轴速度接口，焊缝跟踪上层控制使用前向速度和偏航角速度完成局部视觉跟踪，横向速度通道为后续扩展保留。\\

\section{本章依据说明}

本章依据系统实现资料中的“外部中心点输入下的中心偏差控制方法”“目标丢失与输入超时安全停止方法”“cmd\_vel 到底盘串口协议的执行方法”和本地 NanoOmni 控制板资料。全向底盘车体速度层运动学公式为理论建模表达；四轮目标转速分配、编码器测速、电机 PID 和控制周期依据本地控制板资料及其配套程序抽象得到，应理解为底盘执行层能力，不写成当前上层焊缝跟踪控制已经使用 $v_y$ 横向纠偏或读取里程计闭环。角速度、积分项、速度约束、超时停止和参数表依据系统中已确认的控制逻辑抽象得到。\\

\chapter*{第四章 实机验证与仿真分析}

前两章分别讨论了焊缝目标视觉感知和基于视觉偏差的控制方法。本章在方法分析基础上，围绕系统实现环境、视觉识别功能验证、闭环运行验证和仿真支撑进行说明。由于目前尚未形成完整定量实验结果，本章不写成精度验证章节，而采用“验证目的—验证内容—可观察现象—当前不足”的方式进行分析。\\

\section{4.1 系统实现环境与总体流程}

本文系统面向 ROS1/catkin 环境组织，整体流程由图像输入、视觉识别、目标位置表征、偏差控制、速度输出和执行支撑组成。图像输入提供相机图像或测试视频帧；视觉识别部分完成目标检测与检测框筛选；位置表征部分提取检测框中心并构建有效标志；控制部分根据目标中心与图像中心之间的偏差生成速度指令；执行支撑部分可由底盘接口或 Gazebo 仿真模型消费速度指令。\\

从验证角度看，系统需要依次确认以下内容。第一，输入图像是否能够稳定进入视觉识别流程。第二，目标检测结果是否能够形成可视化检测框。第三，检测框中心和图像中心之间的偏差是否能够输出给控制逻辑。第四，控制逻辑是否能够根据偏差生成速度指令。第五，当检测无效或输入超时时，系统是否能够输出零速度。第六，仿真环境或实机接口是否能够接收速度指令并产生运动响应。\\

系统已具备上述验证入口和功能支撑，但尚不能证明所有验证已经完成并获得定量结果。因此，本章重点说明系统可以如何验证，以及现有实现能支持哪些结论。对于尚未提供截图、日志或数据的部分，本章统一标注为需要补充。\\

\section{4.2 视觉识别功能验证}

视觉识别功能验证的目标是确认系统能够从输入图像中识别焊缝目标区域，并生成后续控制所需的目标中心位置。验证过程不应只关注是否有图像显示，而应关注检测框、目标中心、图像中心线和检测有效状态是否能够形成完整视觉输出。\\

在测试输入条件下，可以使用本地视频或图片作为图像来源，观察检测结果图是否出现焊缝目标检测框。若检测框能够稳定覆盖焊缝目标区域，则说明目标检测流程能够输出可用于位置表征的区域信息。随后，应观察检测框中心是否与图像中心线形成直观偏差关系。该偏差关系是第三章控制方法的输入基础。\\

当前系统可生成包含检测框、目标中心、图像中心线和置信度文字的结果图。该结果图适合放入论文中作为视觉识别功能验证图。但是，目前尚未提供具体截图，因此只能写“系统具备结果图输出功能”，不能写“实验图显示识别效果良好”或“检测准确率达到某数值”。\\

正式论文中，本节建议补充三类图像：第一类为原始输入图像，用于说明视觉输入场景；第二类为检测框可视化图，用于说明目标检测输出；第三类为中心点和图像中心线示意图，用于说明位置偏差构建。若能够补充连续帧结果，还可进一步分析检测框是否稳定、中心点是否连续变化以及是否存在漏检。\\

\section{4.3 焊缝跟踪闭环运行验证}

焊缝跟踪闭环运行验证的目标是确认视觉输出能否驱动控制输出。系统在检测到有效焊缝目标后，将目标中心和图像宽度传递给控制部分。控制部分根据目标中心与图像中心之间的偏差生成速度指令。当目标位于图像中心附近时，系统应输出前进速度；当目标偏离图像中心时，系统应输出转向角速度；当目标无效或输入超时时，系统应输出零速度。\\

闭环验证可分为三个层次。第一是软件链路验证，即不接入真实底盘，只观察视觉输出和速度输出是否随图像变化而变化。该层次能够验证感知—控制数据链路是否成立。第二是仿真验证，即在 Gazebo 环境中观察机器人模型是否能够接收速度指令并产生运动响应。该层次能够验证速度指令与仿真执行接口是否连接。第三是实机验证，即在安全条件下接入真实底盘，观察速度指令是否能够驱动车体沿焊缝目标方向运动。目前系统实现可支撑前两类验证和第三类接口说明，实机闭环效果仍需要补充运行记录。\\

为了使闭环验证具有论文说服力，应记录目标中心输出、速度输出和运行画面之间的对应关系。例如，当焊缝目标偏向图像一侧时，速度输出中的角速度应出现相应变化；当目标重新接近图像中心时，角速度应减小；当目标丢失时，速度输出应变为零。若能够形成图像帧、偏差曲线和速度曲线的对应关系，则第四章可以从现象分析进一步提升为结果分析。\\

目前尚未提供中心偏差曲线、速度输出截图或运行视频，因此本文不能写“闭环跟踪稳定”“实机跟踪误差较小”等结论。可写的结论应限定为系统具备闭环运行验证条件，且已形成由视觉偏差到速度输出的链路。\\

\section{4.4 Gazebo 仿真场景与模型支撑分析}

Gazebo 仿真环境能够为机器人系统提供模型、传感器和运动执行支撑。本文系统包含机器人描述文件、Gazebo 平面运动插件、RGB 相机插件和焊缝相关场景文件。其中，相机插件可为视觉检测提供图像输入，平面运动插件可接收速度指令并驱动机器人模型运动。因此，仿真环境具备支撑感知—控制闭环调试的基础。\\

从论文分析角度看，仿真环境的作用主要体现在三方面。第一，它提供了可重复的焊缝场景，使视觉检测方法能够在固定环境中观察目标识别效果。第二，它提供了机器人模型和相机视角，使图像输入与机器人运动之间形成联系。第三，它提供了速度执行接口，使控制输出能够作用于仿真模型，从而验证感知—控制链路是否能够闭合。\\

系统中设置了曲线焊缝视觉模型和纹理化焊缝场景文件。曲线焊缝场景可用于观察目标中心随路径变化的情况，纹理化场景可用于观察背景变化对检测结果的影响。需要说明的是，纹理资源在跨机器运行前需要确认路径可用。因此，正式论文若要展示纹理化场景，应补充实际运行截图并说明资源路径已经配置完成。\\

对于“半宽焊缝”“正常宽度焊缝”或“不同焊缝宽度对比”等内容，目前尚未形成可直接用于论文的结果截图。因此，本章不能写成已经完成不同宽度焊缝实验。若后续提供正常宽度、半宽或不同背景场景的截图和视频，可在本节增加对比分析。\\

\section{4.5 不同焊缝场景下的运行现象分析}

在缺少定量结果的情况下，不同焊缝场景分析应采用现象分析方式。对于曲线焊缝场景，可观察检测框是否能够随焊缝目标位置变化而变化，目标中心是否连续移动，速度指令是否随偏差变化而调整。对于纹理背景场景，可观察背景纹理是否导致误检、漏检或检测框抖动。对于目标丢失场景，可观察系统是否及时输出零速度。\\

不同场景分析的重点不应放在未提供数据的“精度”上，而应放在系统现象是否符合方法逻辑。例如，当目标位于图像中心附近时，系统应表现为较平稳前进；当目标偏离图像中心时，系统应产生转向控制；当目标检测失败时，系统应停止。若这些现象能够通过截图、视频或日志得到证明，则可以说明系统方法具有基本可行性。\\

【需要用户补充：曲线焊缝场景检测截图、纹理焊缝场景检测截图、不同宽度焊缝对比截图、Gazebo 运行录屏、目标丢失停车截图、中心偏差和速度输出日志。】\\

\section{4.6 实验结果与不足}

从系统实现角度看，当前系统已经具备运行验证入口和功能支撑。系统可以接收图像输入，可以调用 YOLOv5 推理得到检测框，可以提取检测框中心，可以根据中心偏差输出速度指令，可以在无检测或超时时输出零速度，也具备底盘接口和 Gazebo 仿真支撑。\\

当前不能确认的是定量实验结果。目前尚未提供识别准确率、mAP、FPS、跟踪误差、速度响应曲线、丢检率、实机运行日志和严格对比实验。因此，本章目前只能说明系统方法具有可行性验证基础，不能说明系统已经达到某个性能指标。\\

正式论文定稿前，应补充以下材料：检测结果截图、中心偏差输出、速度输出、仿真运行截图、实机运行照片、关键帧图像、运行日志、检测有效帧统计、跟踪偏差统计和安全停止验证。补充这些材料后，本章可以从“可行性验证”扩展为“实验结果分析”。\\

\section{4.7 本章小结}

本章围绕系统验证与仿真分析进行了说明。本文系统已经具备图像输入、视觉检测、中心位置输出、速度控制、零速度停止、底盘接口和 Gazebo 仿真支撑等功能基础。Gazebo 环境中的机器人模型、相机插件、平面运动插件和焊缝场景可用于后续闭环调试。\\

由于目前尚未形成完整定量实验结果，本章没有写成精度验证，而是以可行性验证和现象分析为主。后续若补充截图、视频、运行日志和统计表，可进一步完善视觉识别实验、闭环跟踪实验和不同焊缝场景对比分析。\\

\section{本章依据说明}

本章依据系统实现资料中的“Gazebo 与机器人模型仿真支撑方法”“cmd\_vel 到底盘串口协议的执行方法”“运行流程与关键命令”“实验、结果与评价指标”和“待用户确认的问题”。正文中没有展开具体启动命令，命令和路径只作为事实依据写入 `09\_论文事实依据索引.md`。\\

\chapter*{第五章 总结与展望}

\section{5.1 全文工作总结}

本文围绕爬壁机器人船体表面焊缝跟踪任务，形成了以视觉检测、目标位置表征和视觉偏差控制为核心的研究内容。全文按照“视觉感知—目标位置表征—偏差构建—运动控制—闭环跟踪—实机/仿真验证”的技术路线展开，重点讨论了如何从输入图像中获得焊缝目标位置，并将该位置转化为机器人运动控制量。\\

在视觉感知方面，本文采用基于 YOLOv5 推理的目标检测方法识别焊缝区域。系统对输入图像进行预处理后，通过目标检测模型获得焊缝目标检测框，并经过候选框筛选和坐标恢复得到原图坐标系下的目标区域。由于目前尚未整理形成训练数据集、类别名称和检测精度等完整材料，因此本文没有对模型训练过程和检测指标作确定性结论。\\

在目标位置表征方面，本文采用检测框中心作为焊缝目标位置的几何表达。通过计算检测框左右边界均值获得目标横向中心，并以图像宽度的一半作为图像中心参考，建立目标中心与图像中心之间的横向偏差。该方法将视觉检测输出转化为控制系统可使用的简洁输入，是感知与控制连接的关键。\\

在运动控制方面，本文围绕基于视觉偏差的速度调节方法和面向视觉闭环的运动优化策略展开。系统结合麦克纳姆轮全向底盘的速度接口建立车体速度层运动学表达，根据目标中心相对图像中心的偏差生成角速度，并根据偏差和转向强度约束线速度。当目标接近图像中心时，机器人保持前进；当目标偏离图像中心时，机器人进行转向修正；当检测结果无效或输入超时时，系统输出零速度并清零积分状态。该控制方法通过归一化误差、小增益积分修正、积分分离、积分限幅和线速度-角速度耦合约束，提高了视觉闭环跟踪的偏差修正能力和运动平稳性。\\

在系统验证方面，本文整理了视觉识别功能验证、闭环运行验证和 Gazebo 仿真支撑分析。由于缺少完整实验截图、日志和数据表，第四章采用可行性验证和现象分析方式，没有给出尚未被实验材料支撑的精度、速度和误差结论。\\

\section{5.2 方法与系统实现总结}

本文方法的核心在于使用检测框中心完成焊缝目标位置表征。该方法没有依赖语义分割、中心线拟合或三维重建，而是围绕当前系统所需的横向偏差信息建立简洁的控制输入。虽然该表征方式相对简单，但它能够直接提供控制所需的横向位置信息，适合构建视觉偏差控制链路。\\

系统实现方面，本文基于 ROS1/catkin 组织焊缝跟踪系统，包含图像输入、视觉检测、位置表征、偏差控制、执行接口和 Gazebo 仿真支撑。实机路径可通过底盘接口消费速度指令，且底盘接口支持前向速度、横向速度和偏航角速度三个分量；仿真路径可通过 Gazebo 模型和插件支撑闭环调试。该实现方式使感知算法、控制逻辑和底层执行接口相对分离，为后续替换图像来源、调整检测权重和继续扩展横向速度纠偏提供了基础。\\

需要强调的是，当前代码材料中的麦克纳姆轮平台主要承担移动执行平台作用，论文工作的重点是面向船体表面焊缝作业背景构建任务层感知与控制链路，而不是展开爬壁结构或吸附机构设计。本文将视觉检测输出进一步组织为目标中心、图像宽度和有效标志，将该位置表征转换为归一化偏差，并据此生成角速度、线速度和停止状态。通过这一过程，移动平台不再只是执行外部给定速度，而是能够依据视觉偏差形成面向焊缝目标的局部闭环响应。\\

从论文方法角度看，本文形成了较完整的处理逻辑：首先通过目标检测获得焊缝区域；然后通过检测框中心获得目标位置；再通过图像中心参考构建横向偏差；最后根据偏差生成速度控制量。该逻辑使视觉信息能够逐步转化为运动控制输入，而不是停留在单纯检测或单纯控制层面。\\

\section{5.3 当前工作的不足}

当前工作的不足主要体现在以下方面。\\

第一，模型训练和评价信息不足。目前系统中存在可供推理使用的模型权重和 YOLOv5 工具链，但缺少焊缝数据集配置、类别名称、训练日志、损失曲线、mAP 和 FPS 等指标。因此，论文无法给出检测模型性能的定量结论。\\

第二，视觉位置表征较为简单。系统使用检测框横向中心作为控制输入，没有进一步提取焊缝中心线、焊缝宽度、目标曲率或连续轨迹。当焊缝形态复杂、目标遮挡或检测框抖动时，仅依赖检测框中心可能会影响控制稳定性。\\

第三，运动优化策略仍属于视觉偏差闭环下的参数化控制优化。当前控制逻辑已经在比例偏差基础上加入小增益积分项、积分分离、积分限幅和速度耦合约束，但没有引入微分项、里程计反馈、动力学模型或轨迹预测，也未实现 MPC、Stanley 或滑模控制等方法。虽然底盘接口支持麦克纳姆轮全向移动所需的横向速度通道，但当前焊缝跟踪上层控制仍保持 $v_y=0$，未实现基于横向平移的纠偏控制。因此，在复杂路径、高速运动或视觉噪声较大场景下，仍需要进一步调参和优化。\\

第四，实验材料不足。目前尚未形成完整实机闭环视频、Gazebo 定量仿真结果、识别成功率、跟踪误差、丢检率、速度响应曲线和对比实验。第四章只能写可行性验证，不能形成充分的结果分析。\\

\section{5.4 后续研究展望}

后续工作首先应补充实机场景数据。可以通过真实相机采集船体表面或模拟船体表面的焊缝图像，或从机器人运行视频中抽取关键帧，建立用于训练、验证和测试的数据集，并明确类别定义、标注规则和数据划分方式。在此基础上，可以针对仿真场景和实机场景分别训练或微调权重，提高检测模型在不同光照、不同背景和不同视角下的鲁棒性。\\

其次，应完善实验评价指标。建议增加检测有效帧比例、识别置信度变化、丢检率、中心偏差均值、最大偏差、速度输出范围和安全停止响应时间等指标。通过这些指标，可以将当前可行性验证扩展为较完整的实验结果分析。\\

再次，可进一步优化控制参数和全向运动利用方式。当前控制策略适合作为基础实现，后续可根据实机运行数据调整角速度增益、积分增益、线速度阈值、积分分离阈值和停止条件。当前底盘资料已能说明下位机侧具备三轴速度解析、四轮目标转速分配和电机速度闭环能力；后续若要利用 $v_y$ 横向速度进行焊缝偏差修正，还需要完成实车固件一致性确认、横向速度通道测试、状态反馈记录和上层控制策略设计。若未来采用 MPC、Stanley 或其他控制方法，应补充相应实现方案和实验结果，不能仅在论文中进行理论描述。\\

最后，应提高仿真环境与实机场景的一致性。可完善 Gazebo 焊缝场景、纹理资源路径和机器人模型配置，补充正常宽度与半宽焊缝等不同场景验证材料，并将仿真现象与实机运行现象进行对照分析。\\

\section{本章小结}

本章总结了全文围绕焊缝目标检测、目标中心表征、视觉偏差控制、运动优化控制和闭环验证所开展的工作。当前系统已经形成较清晰的软件实现链路，但在训练数据、定量评价、实机验证、横向速度纠偏和控制参数优化方面仍需补充。后续研究应重点围绕真实数据采集、模型评估、控制参数优化、全向运动能力利用和实验材料完善展开。\\

\section{本章依据说明}

本章依据系统实现资料中的“项目总体定位”“整体技术路线”“实验、结果与评价指标”“明确不能写的内容”和“待用户确认的问题”。展望内容均作为未来工作提出，没有写成当前系统已经完成的功能。\\

\end{document}
